% 数学ノート集・共通テンプレート v2（積分論2026年度第1部の版面を参考）
% A5判。章内で定義・定理・例・注意等の番号を共通化。
% コンパイル： platex -interaction=nonstopmode math_note_template_v2.tex (2回)
%              dvipdfmx math_note_template_v2.dvi
\documentclass[a5paper,10pt,onecolumn]{jsarticle}
\usepackage[dvipdfmx,top=22mm,bottom=20mm,inner=19mm,outer=18mm,headheight=13pt,headsep=6mm]{geometry}
\usepackage{amsmath,amssymb,amsthm,mathtools}
\usepackage{enumitem,fancyhdr,titlesec}
\usepackage[dvipdfmx,bookmarks=true,colorlinks=false]{hyperref}
\usepackage{pxjahyper}

% ----- 差し替える文書情報 -----
\newcommand{\NoteTitle}{フィルター・超フィルターと位相的収束}
\newcommand{\NoteSeries}{数学ノート集}

% ----- 文字・組版 -----
% 参考PDFの質素な明朝＋Computer Modernの印象を引き継ぎながら，
% A5判で窮屈にならないよう本文の行間を少し確保する．
\linespread{1.13}
\setlength{\parindent}{1zw}
\setlength{\parskip}{0pt}
\setlength{\abovedisplayskip}{6pt plus 2pt minus 2pt}
\setlength{\belowdisplayskip}{6pt plus 2pt minus 2pt}
\setlength{\abovedisplayshortskip}{3pt plus 2pt minus 1pt}
\setlength{\belowdisplayshortskip}{4pt plus 2pt minus 1pt}
\setlist[enumerate]{topsep=3pt,itemsep=1pt,parsep=0pt,leftmargin=2.4zw}
\titleformat{\section}{\normalfont\normalsize\bfseries}{\thesection}{0.8zw}{}
\titlespacing*{\section}{0pt}{2.0ex plus .4ex}{.9ex}
\titleformat{\subsection}{\normalfont\normalsize\bfseries}{\thesubsection}{0.7zw}{}
\titlespacing*{\subsection}{0pt}{1.5ex plus .3ex}{.6ex}
\raggedbottom

% ----- 柱とページ番号：ページ上部、罫線なし -----
\fancypagestyle{mathnote}{%
  \fancyhf{}%
  \fancyhead[L]{\footnotesize \NoteSeries\ （\NoteTitle）}%
  \fancyhead[R]{\footnotesize \thepage}%
  \renewcommand{\headrulewidth}{0pt}%
  \renewcommand{\footrulewidth}{0pt}%
}
\pagestyle{mathnote}

% ----- 節内通し番号の定義・定理群 -----
\newtheoremstyle{mathnotetheorem}
 {4.5pt}{4.5pt}{\normalfont}{}%
 {\bfseries}{}{0.65zw}%
 {\thmname{#1}\thmnumber{ #2}\thmnote{（#3）}}
\theoremstyle{mathnotetheorem}
\newtheorem{theorem}{定理}[section]
\newtheorem{definition}[theorem]{定義}
\newtheorem{lemma}[theorem]{補題}
\newtheorem{proposition}[theorem]{命題}
\newtheorem{corollary}[theorem]{系}
\newtheorem{example}[theorem]{例}
\newtheorem{remark}[theorem]{注意}
\renewcommand{\proofname}{証明}
\renewcommand{\qedsymbol}{$\square$}
\makeatletter
\renewenvironment{proof}[1][\proofname]{%
  \par\pushQED{\qed}\normalfont\topsep4\p@\@plus2\p@\relax
  \trivlist\item[\hskip\labelsep\bfseries #1]\ignorespaces
}{\popQED\endtrivlist\@endpefalse}
\makeatother

% ----- 使うノートで適宜拡張する記号 -----
\newcommand{\R}{\mathbb{R}}
\newcommand{\norm}[1]{\lVert #1\rVert}
\newcommand{\abs}[1]{\lvert #1\rvert}

\begin{document}
\section*{読了後にできること}
本稿を読むための前提は，フィルター，フィルターの収束，Hausdorff性．
本稿の具体的な到達課題は，フィルターの収束先の一意性を証明する．

\section{フィルターとフィルター基}

本稿では集合論の通常の公理系を前提とし，
超フィルターの存在を示す箇所でZornの補題を用いる．
位相空間は特記しない限りHausdorffとは限らない．

\begin{definition}
空でない集合 $X$ 上の部分集合族
$\mathcal F\subset\mathcal P(X)$ が\emph{フィルター}であるとは，
次の三条件を満たすことをいう．
\begin{enumerate}
\item $\varnothing\notin\mathcal F$．
\item $A,B\in\mathcal F$ なら $A\cap B\in\mathcal F$．
\item $A\in\mathcal F$，$A\subset B\subset X$ なら
$B\in\mathcal F$．
\end{enumerate}
\end{definition}

\begin{example}
$x\in X$ に対する
$\mathcal U_x=\{A\subset X:x\in A\}$ はフィルターである．
$X$ が無限集合なら
$\mathcal F_{\mathrm{cof}}
 =\{A\subset X:X\setminus A\text{ は有限}\}$
もフィルターである．
\end{example}
\begin{proof}
前者は，$x$ を含む集合の共通部分も $x$ を含むことから
各公理を満たす．
後者では二つの有限集合の和集合が有限であり，
無限集合 $X$ の補有限部分集合は空でないので
各公理を満たす．
\end{proof}

\begin{definition}
$\mathcal B\subset\mathcal P(X)$ が
\emph{フィルター基}であるとは，
$\mathcal B$ が空でなく，空集合を含まず，
任意の $B_1,B_2\in\mathcal B$ に対し
$B_3\in\mathcal B$ で
$B_3\subset B_1\cap B_2$ となるものが存在することをいう．
\end{definition}

\begin{proposition}[生成されるフィルター]\label{prop:generated}
フィルター基 $\mathcal B$ に対して
\[
 \langle\mathcal B\rangle
 =\{A\subset X:\text{ある }B\in\mathcal B\text{ が }B\subset A\}
\]
はフィルターである．
これは $\mathcal B$ を含むフィルターのうち最小のものである．
\end{proposition}
\begin{proof}
$\mathcal B$ の元は非空なので
$\varnothing\notin\langle\mathcal B\rangle$ である．
$A_1,A_2$ がそれぞれ $B_1,B_2\in\mathcal B$ を含むなら，
基の公理によって
$B_3\subset B_1\cap B_2\subset A_1\cap A_2$
となる $B_3\in\mathcal B$ があり，
$A_1\cap A_2\in\langle\mathcal B\rangle$ となる．
上方閉性は定義から従う．
$\mathcal B\subset\langle\mathcal B\rangle$ であり，
$\mathcal B$ を含む任意のフィルターは上方閉性から
$\langle\mathcal B\rangle$ を含む．
\end{proof}

\begin{proposition}[有限交叉性]\label{prop:fip}
集合族 $\mathcal A\subset\mathcal P(X)$ を含む
フィルターが存在するための必要十分条件は，
$\mathcal A$ の任意の有限個の元の共通部分が
非空であることである．
この条件を\emph{有限交叉性}という．
\end{proposition}
\begin{proof}
フィルターが $\mathcal A$ を含むなら，
有限共通部分もフィルターに属するため空でない．
逆に有限交叉性が成立するなら，
$\mathcal A$ の元の有限個の共通部分全体を
$\mathcal B$ とする．
空個の共通部分を $X$ として含めれば
$\mathcal B$ はフィルター基となる．
命題\ref{prop:generated}の生成フィルターは
$\mathcal A$ を含む．
\end{proof}

\section{フィルターによる収束}

\begin{definition}
$X$ を位相空間，$\mathcal F$ を $X$ 上のフィルターとする．
$x\in X$ の近傍全体を $\mathcal N(x)$ と書く．
\[
 \mathcal F\longrightarrow x
 \quad\Longleftrightarrow\quad
 \mathcal N(x)\subset\mathcal F
\]
によってフィルターの\emph{収束}を定める．
また $x$ が $\mathcal F$ の\emph{集積点}であるとは，
任意の $A\in\mathcal F$ と $x$ の任意の近傍 $U$ に対し
$A\cap U\ne\varnothing$ となることをいう．
\end{definition}

\begin{proposition}[集積点の特徴づけ]\label{prop:cluster}
$x$ がフィルター $\mathcal F$ の集積点であることと
\[
 x\in\bigcap_{A\in\mathcal F}\overline A
\]
が成り立つことは同値である．
また $\mathcal F\to x$ なら $x$ は集積点である．
\end{proposition}
\begin{proof}
$x\in\overline A$ であることは
$x$ の全ての近傍 $U$ が $A$ と交わることと同値である．
全ての $A\in\mathcal F$ に対してこの条件を課せば
前半が従う．
$\mathcal F\to x$ なら $U\in\mathcal F$ なので
$A\cap U\in\mathcal F$ であり，
フィルターは空集合を含まないため
$A\cap U\ne\varnothing$ となる．
\end{proof}

\begin{proposition}[連続写像とフィルター]\label{prop:continuous}
$f:X\to Y$ を写像とし，
$\mathcal F$ の\emph{像フィルター}を
\[
 f_*\mathcal F
 =\{B\subset Y:f^{-1}(B)\in\mathcal F\}
\]
と定める．これは $Y$ 上のフィルターである．
$f$ が連続なら
$\mathcal F\to x$ から
$f_*\mathcal F\to f(x)$ が従う．
逆に任意の $x$ について
$\mathcal N(x)$ の像が $f(x)$ に収束するなら
$f$ は連続である．
\end{proposition}
\begin{proof}
$f^{-1}(\varnothing)=\varnothing$，
逆像は有限共通部分と包含を保つから
$f_*\mathcal F$ はフィルターである．
$f$ が連続で $\mathcal F\to x$ とする．
$f(x)$ の近傍 $V$ について $f^{-1}(V)$ は
$x$ の近傍なので $\mathcal F$ に属する．
従って $V\in f_*\mathcal F$ となり収束する．
逆に $\mathcal N(x)$ の像が $f(x)$ に収束するなら，
任意の $f(x)$ の近傍 $V$ について
$f^{-1}(V)\in\mathcal N(x)$ である．
従って $f$ は全ての $x$ で連続である．
\end{proof}

\begin{theorem}[Hausdorff性と極限の一意性]\label{thm:hausdorff}
位相空間 $X$ がHausdorffであることと，
全てのフィルターの収束先が高々一つであることは同値である．
\end{theorem}
\begin{proof}
$X$ がHausdorffで
$\mathcal F\to x$ かつ $\mathcal F\to y$ とする．
$x\ne y$ なら交わらない開近傍 $U,V$ が存在する．
両方とも $\mathcal F$ に属するので
$U\cap V=\varnothing$ が $\mathcal F$ に属し矛盾する．

逆に $X$ がHausdorffでないとする．
異なる二点 $x,y$ で，
任意の近傍 $U\ni x$，$V\ni y$ が交わるものが存在する．
集合族
$\{U\cap V:U\in\mathcal N(x),V\in\mathcal N(y)\}$
はフィルター基である．
そこから生成されるフィルターは
$\mathcal N(x)$ と $\mathcal N(y)$ をともに含むため
$x,y$ の両方に収束する．
従って極限の一意性を満たさない．
\end{proof}

\section{ネットとフィルターの相互変換}

\begin{definition}
有向集合 $D$ 上の写像 $x:D\to X$ を
\emph{ネット}といい $(x_d)_{d\in D}$ と書く．
ネットが $x\in X$ に収束するとは，
$x$ の任意の近傍 $U$ について
ある $d_0\in D$ が存在して
$d\ge d_0$ なら $x_d\in U$ となることをいう．
\end{definition}

\begin{proposition}[ネットが定めるフィルター]\label{prop:net-filter}
ネット $(x_d)_{d\in D}$ に対し
\[
 \mathcal F_x=
 \{A\subset X:\text{ある }d_0\in D
 \text{ に対し }d\ge d_0\Rightarrow x_d\in A\}
\]
と置くと $\mathcal F_x$ はフィルターである．
また $x_d\to x$ と $\mathcal F_x\to x$ は同値である．
\end{proposition}
\begin{proof}
尾部の像
$T_{d_0}=\{x_d:d\ge d_0\}$ は非空である．
二つの尾部 $T_{d_1},T_{d_2}$ に対して
有向性から $d_3\ge d_1,d_2$ を取れば
$T_{d_3}\subset T_{d_1}\cap T_{d_2}$ となる．
従って尾部の像の族はフィルター基であり，
$\mathcal F_x$ はその生成フィルターである．
$x_d\to x$ とは，
任意の $x$ の近傍 $U$ が十分先の尾部を含むことであり，
これは $\mathcal N(x)\subset\mathcal F_x$ と同値である．
\end{proof}

\begin{theorem}[フィルターからネットへ]\label{thm:filter-net}
$X$ 上の任意のフィルター $\mathcal F$ に対して，
あるネット $(x_d)$ が存在して
\[
 \mathcal F_x=\mathcal F
\]
が成立する．
従ってフィルター収束とネット収束は同値な概念である．
\end{theorem}
\begin{proof}
\[
 D=\{(A,a):A\in\mathcal F,\ a\in A\}
\]
と置く．
$(A,a)\preceq(B,b)$ を $B\subset A$ によって定義する．
これは反射的・推移的な前順序であり，
有向集合として利用できる．
実際二つの元 $(A,a),(B,b)$ に対し
$A\cap B\in\mathcal F$ は非空なので
$c\in A\cap B$ を取れば，
$(A\cap B,c)$ は両者の上界である．
ネットを $x_{(A,a)}=a$ と定める．

任意の $A\in\mathcal F$ と $a\in A$ に対し，
$(A,a)\preceq(B,b)$ なら $b\in B\subset A$ なので
ネットは最終的に $A$ に入る．
従って $\mathcal F\subset\mathcal F_x$ である．
逆に $S\in\mathcal F_x$ なら，
ある $(A,a)\in D$ より先の全てのネット値が
$S$ に属する．
任意の $b\in A$ に対し $(A,a)\preceq(A,b)$ だから
$b\in S$ となり，$A\subset S$ を得る．
$A\in\mathcal F$ の上方閉性から
$S\in\mathcal F$ である．従って等号が成り立つ．
\end{proof}

\section{超フィルターと選択原理}

\begin{definition}
フィルター $\mathcal U$ が\emph{超フィルター}であるとは，
包含関係について極大な真のフィルターであることをいう．
\end{definition}

\begin{lemma}[超フィルターの二値性]\label{lem:ultra-character}
フィルター $\mathcal U$ について，次は同値である．
\begin{enumerate}
\item $\mathcal U$ は超フィルターである．
\item 任意の $A\subset X$ に対して
$A\in\mathcal U$ または
$X\setminus A\in\mathcal U$ のいずれか一方が成り立つ．
\end{enumerate}
\end{lemma}
\begin{proof}
(1) を仮定する．
$A\notin\mathcal U$ とする．
もし全ての $B\in\mathcal U$ について
$B\cap A\ne\varnothing$ なら，
$\mathcal U\cup\{A\}$ は有限交叉性を持つ．
命題\ref{prop:fip}により
$\mathcal U$ を真に含むフィルターが存在し，
極大性に反する．
従って $B\in\mathcal U$ で
$B\cap A=\varnothing$ となるものが存在する．
すると $B\subset X\setminus A$ なので
$X\setminus A\in\mathcal U$ である．
両方が同時に属すれば共通部分が空となるので，
ちょうど一方が属する．

(2) を仮定し，
$\mathcal U$ を真に含むフィルター $\mathcal V$ が
存在するとする．
$A\in\mathcal V\setminus\mathcal U$ を取ると
$X\setminus A\in\mathcal U\subset\mathcal V$ である．
従って空集合が $\mathcal V$ に属し矛盾する．
ゆえに $\mathcal U$ は極大である．
\end{proof}

\begin{theorem}[超フィルターの延長]\label{thm:ultra-extension}
集合 $X$ 上の任意のフィルター $\mathcal F$ は，
ある超フィルター $\mathcal U$ に含まれる．
\end{theorem}
\begin{proof}
$\mathcal F$ を含むフィルター全体を
包含関係で順序づけた集合 $\mathcal P$ を考える．
$\mathcal P$ は $\mathcal F$ を含むので空でない．
任意の鎖 $\mathcal C\subset\mathcal P$ に対して
$\mathcal V=\bigcup_{\mathcal G\in\mathcal C}\mathcal G$
と置く．
空集合はどの $\mathcal G$ にも属さないから
$\mathcal V$ にも属さない．
$A,B\in\mathcal V$ なら
$A\in\mathcal G_1$，$B\in\mathcal G_2$ となる
$\mathcal G_1,\mathcal G_2\in\mathcal C$ が存在する．
鎖であるから一方が他方を含み，
従って $A,B$ は一つのフィルターに属する．
ゆえに $A\cap B\in\mathcal V$ である．
上方閉性もいずれかの $\mathcal G$ で成り立つ．
従って $\mathcal V$ は $\mathcal P$ に属し，
鎖の上界となる．
Zornの補題を適用すると $\mathcal P$ は極大元を持ち，
それは $\mathcal F$ を含む超フィルターである．
\end{proof}

\begin{corollary}[自由超フィルター]\label{cor:free-ultra}
$X$ が無限集合なら，有限集合を一つも含まない
超フィルターが $X$ 上に存在する．
\end{corollary}
\begin{proof}
補有限フィルター $\mathcal F_{\mathrm{cof}}$ を
超フィルター $\mathcal U$ へ延長する．
もし有限集合 $A$ が $\mathcal U$ に属すれば
その補集合 $X\setminus A$ は
$\mathcal F_{\mathrm{cof}}\subset\mathcal U$ に属し，
空集合が $\mathcal U$ に属することになり矛盾する．
\end{proof}

\section{コンパクト性の特徴づけ}

\begin{lemma}[超フィルターでは集積と収束が一致]
\label{lem:ultra-cluster}
超フィルター $\mathcal U$ の集積点 $x$ は
$\mathcal U$ の収束先である．
\end{lemma}
\begin{proof}
$x$ が集積点なら任意の $x$ の近傍 $V$ と
任意の $A\in\mathcal U$ について
$A\cap V\ne\varnothing$ である．
もし $V\notin\mathcal U$ なら，
補題\ref{lem:ultra-character}から
$X\setminus V\in\mathcal U$ となる．
この $A=X\setminus V$ は $V$ と交わらないので矛盾する．
従って全ての近傍 $V$ が $\mathcal U$ に属し，
$\mathcal U\to x$ である．
\end{proof}

\begin{theorem}[コンパクト性の超フィルター判定]
\label{thm:compact-ultra}
位相空間 $X$ について，次は同値である．
\begin{enumerate}
\item $X$ はコンパクトである．
\item $X$ 上の全ての超フィルターは $X$ 内に収束する．
\item $X$ 上の全てのフィルターは集積点を持つ．
\end{enumerate}
\end{theorem}
\begin{proof}
(1) から (3) を示す．
$\mathcal F$ を任意のフィルターとする．
各 $A\in\mathcal F$ の閉包 $\overline A$ を考えると，
\[
 \bigcap_{j=1}^n\overline{A_j}
 \supset\overline{\bigcap_{j=1}^nA_j}
\]
である．
$\bigcap_jA_j\in\mathcal F$ は非空なので，
これらの閉包は有限交叉性を持つ．
コンパクト性によって
$\bigcap_{A\in\mathcal F}\overline A\ne\varnothing$
となり，命題\ref{prop:cluster}から集積点が存在する．

(3) から (2) を示す．
超フィルター $\mathcal U$ はフィルターなので
集積点 $x$ を持つ．補題\ref{lem:ultra-cluster}から
$\mathcal U\to x$ である．

(2) から (1) を示す．
$X$ がコンパクトでないと仮定する．
有限部分被覆を持たない開被覆 $\{U_i:i\in I\}$ を取る．
閉集合 $F_i=X\setminus U_i$ の族は有限交叉性を持つから，
命題\ref{prop:fip}と定理\ref{thm:ultra-extension}により，
全ての $F_i$ を含む超フィルター $\mathcal U$ が存在する．
(2) により $\mathcal U\to x$ となる $x\in X$ を取る．
開被覆なので $x\in U_{i_0}$ となる $i_0$ が存在する．
$U_{i_0}$ は $x$ の近傍だから $\mathcal U$ に属する．
一方 $F_{i_0}\in\mathcal U$ なので
空集合が $\mathcal U$ に属し矛盾する．
よって $X$ はコンパクトである．
\end{proof}

\begin{corollary}[集積するネット]\label{cor:compact-net}
$X$ がコンパクトであることと，
任意のネットが集積点を持つことは同値である．
\end{corollary}
\begin{proof}
ネットの尾部フィルターを命題\ref{prop:net-filter}で作れば，
フィルターの集積点とネットの集積点
（任意の尾部において任意の近傍へ入る点が存在すること）
は同値である．
定理\ref{thm:filter-net}により任意のフィルターは
ネットの尾部フィルターとして実現される．
従って定理\ref{thm:compact-ultra}の (3) に帰着する．
\end{proof}

\section{積位相とTychonoffの定理}

\begin{lemma}[フィルターの積空間収束]\label{lem:product-convergence}
$X=\prod_{i\in I}X_i$ に積位相を入れる．
フィルター $\mathcal F$ が $x=(x_i)_{i\in I}$ に
収束するための必要十分条件は，
任意の座標射影 $\pi_i:X\to X_i$ に対して
$(\pi_i)_*\mathcal F\to x_i$ が成り立つことである．
\end{lemma}
\begin{proof}
必要性は座標射影の連続性と
命題\ref{prop:continuous}から従う．
逆に各座標像が $x_i$ に収束すると仮定する．
$x$ の任意の基本近傍は
\[
 \bigcap_{j=1}^n\pi_{i_j}^{-1}(U_j)
\]
という形であり，$U_j$ は $x_{i_j}$ の近傍である．
仮定から各 $\pi_{i_j}^{-1}(U_j)$ は
$\mathcal F$ に属するので，
有限共通部分も $\mathcal F$ に属する．
従って $x$ の全ての近傍が $\mathcal F$ に属し，
$\mathcal F\to x$ である．
\end{proof}

\begin{theorem}[Tychonoffの積定理]\label{thm:tychonoff}
任意の添字集合 $I$ に対して
各 $X_i$ がコンパクトなら，
積位相を持つ
\[
 X=\prod_{i\in I}X_i
\]
はコンパクトである．
\end{theorem}
\begin{proof}
空の因子 $X_i$ が存在すれば積空間は空となりコンパクトである．
以下，各 $X_i$ は空でないとする．
選択公理によって $X$ は空でない．
$X$ 上の任意の超フィルター $\mathcal U$ を取る．
各座標射影による像
$(\pi_i)_*\mathcal U$ が超フィルターであることを示す．
任意の $A\subset X_i$ について，
$\pi_i^{-1}(A)$ と
$\pi_i^{-1}(X_i\setminus A)$ は補集合の関係にある．
超フィルターの二値性から
そのいずれかちょうど一方が $\mathcal U$ に属する．
従って $(\pi_i)_*\mathcal U$ は超フィルターである．

$X_i$ はコンパクトなので
定理\ref{thm:compact-ultra}から
$(\pi_i)_*\mathcal U\to x_i$ となる
$x_i\in X_i$ が存在する．
選択公理によってそれらを同時に選び
$x=(x_i)_{i\in I}\in X$ とする．
補題\ref{lem:product-convergence}から
$\mathcal U\to x$ である．
従って $X$ 上の全ての超フィルターは収束し，
再び定理\ref{thm:compact-ultra}によって
$X$ はコンパクトである．
\end{proof}

\begin{remark}
Zornの補題を用いた超フィルター延長と，
Tychonoffの積定理の一般形はいずれも選択原理と深く関係する．
ただしフィルター・ネットの収束の定義および
その相互変換は選択原理を持ち出さずに扱える．
本稿の議論は，次の主題である超実数の超冪構成の
準備にもなる．
\end{remark}

\section{到達課題の完全解答}
\noindent\textbf{Hausdorff空間でのフィルター極限．}
\begin{proposition}
Hausdorff空間 $X$ のフィルター $\mathcal F$ は
相異なる二点に同時に収束しない．
\end{proposition}
\begin{proof}
$\mathcal F$ が $x\ne y$ に収束すると仮定する．
Hausdorff性から $x$ と $y$ の互いに素な
開近傍 $U,V$ を選べる．
収束の定義より $U,V\in\mathcal F$ である．
フィルターは有限交叉で閉じているので
$U\cap V=\varnothing$ も $\mathcal F$ に属することになるが，
真のフィルターは空集合を含まないため矛盾する．
\end{proof}
\noindent\textbf{到達確認．}
無限補有限空間では極限は一意か．
\par\noindent\textbf{解答．}
非空開集合全体からなる補有限フィルターは
各点の全近傍を含むので，全ての点に収束する．
従って極限は一意でない．
\end{document}